%
%%%%%%%%%%%%%%%%%%%%%%%%%%%%%%%%%%%%%%%%%%%%%%%%%%%%%%%%%%%%%%%%%%%%%%%%%%%%%%
% The 2nd Trial by Andrea Tino is licensed under CC BY-NC-SA 4.0. To view a   %
% copy of this license, visit                                                 %
% http://creativecommons.org/licenses/by-nc-sa/4.0/                           %
%%%%%%%%%%%%%%%%%%%%%%%%%%%%%%%%%%%%%%%%%%%%%%%%%%%%%%%%%%%%%%%%%%%%%%%%%%%%%%%

\documentclass[twoside,symmetric,justified]{tufte-book}

\hypersetup{colorlinks} % colored hyperlinks (e.g., for onscreen viewing)

%%
% Book metadata
\title{The 2nd Trial}
\author[Andrea Tino]{Andrea Tino}
\publisher{Issue I - Year 2019}

%\usepackage{microtype}

%%
% For nicely typeset tabular material
\usepackage{booktabs}

% TIKZ
\usepackage{tikz}
\usetikzlibrary{plotmarks}
\usetikzlibrary{patterns}
\usetikzlibrary{decorations.markings}
\usetikzlibrary{math}
\usetikzlibrary{matrix}
\usetikzlibrary{arrows,tikzmark,shadows,positioning}
\usepackage{tikzscale}

%%
% For graphics / images
\usepackage{graphicx}
\setkeys{Gin}{width=\linewidth,totalheight=\textheight,keepaspectratio}
\graphicspath{{graphics/}}

% The fancyvrb package lets us customize the formatting of verbatim
% environments.  We use a slightly smaller font.
\usepackage{fancyvrb}
\fvset{fontsize=\normalsize}

%%
% Prints argument within hanging parentheses (i.e., parentheses that take
% up no horizontal space).  Useful in tabular environments.
\newcommand{\hangp}[1]{\makebox[0pt][r]{(}#1\makebox[0pt][l]{)}}

%%
% Prints an asterisk that takes up no horizontal space.
% Useful in tabular environments.
\newcommand{\hangstar}{\makebox[0pt][l]{*}}

%%
% Prints a trailing space in a smart way.
\usepackage{xspace}

\newcommand{\TST}{The 2nd Trial\xspace}

% Prints the month name (e.g., January) and the year (e.g., 2008)
\newcommand{\monthyear}{%
  \ifcase\month\or January\or February\or March\or April\or May\or June\or
  July\or August\or September\or October\or November\or
  December\fi\space\number\year
}


% Prints an epigraph and speaker in sans serif, all-caps type.
\newcommand{\openepigraph}[2]{%
  %\sffamily\fontsize{14}{16}\selectfont
  \begin{fullwidth}
  \sffamily\large
  \begin{doublespace}
  \noindent\allcaps{#1}\\% epigraph
  \noindent\allcaps{#2}% author
  \end{doublespace}
  \end{fullwidth}
}

% Inserts a blank page
\newcommand{\blankpage}{\newpage\hbox{}\thispagestyle{empty}\newpage}

% Creates all needed for a new paragraph with side note marker
% 1. The short caption
% 2. The title
% 3. The ref
\newcommand{\newpar}[4][0pt]{%
    \paragraph[#2]{#3}
    \label{#4}
    \marginnote[#1]{\smallcaps{\nameref{#4}}}
}

\usepackage{units}

% Ensures side boxes are places on the correct side by tracking the page they fall in
\usepackage{mparhack}

% Typesets the font size, leading, and measure in the form of 10/12x26 pc.
\newcommand{\measure}[3]{#1/#2$\times$\unit[#3]{pc}}

% Macros for typesetting the documentation
\newcommand{\hlred}[1]{\textcolor{Maroon}{#1}}% prints in red
\newcommand{\hangleft}[1]{\makebox[0pt][r]{#1}}
\newcommand{\hairsp}{\hspace{1pt}}% hair space
\newcommand{\hquad}{\hskip0.5em\relax}% half quad space
\newcommand{\TODO}{\textcolor{red}{\bf TODO!}\xspace}
\newcommand{\ie}{\textit{i.\hairsp{}e.}\xspace}
\newcommand{\eg}{\textit{e.\hairsp{}g.}\xspace}
\newcommand{\na}{\quad--}% used in tables for N/A cells
\providecommand{\XeLaTeX}{X\lower.5ex\hbox{\kern-0.15em\reflectbox{E}}\kern-0.1em\LaTeX}
\newcommand{\tXeLaTeX}{\XeLaTeX\index{XeLaTeX@\protect\XeLaTeX}}
% \index{\texttt{\textbackslash xyz}@\hangleft{\texttt{\textbackslash}}\texttt{xyz}}
\newcommand{\tuftebs}{\symbol{'134}}% a backslash in tt type in OT1/T1
\newcommand{\doccmdnoindex}[2][]{\texttt{\tuftebs#2}}% command name -- adds backslash automatically (and doesn't add cmd to the index)
\newcommand{\doccmddef}[2][]{%
  \hlred{\texttt{\tuftebs#2}}\label{cmd:#2}%
  \ifthenelse{\isempty{#1}}%
    {% add the command to the index
      \index{#2 command@\protect\hangleft{\texttt{\tuftebs}}\texttt{#2}}% command name
    }%
    {% add the command and package to the index
      \index{#2 command@\protect\hangleft{\texttt{\tuftebs}}\texttt{#2} (\texttt{#1} package)}% command name
      \index{#1 package@\texttt{#1} package}\index{packages!#1@\texttt{#1}}% package name
    }%
}% command name -- adds backslash automatically
\newcommand{\doccmd}[2][]{%
  \texttt{\tuftebs#2}%
  \ifthenelse{\isempty{#1}}%
    {% add the command to the index
      \index{#2 command@\protect\hangleft{\texttt{\tuftebs}}\texttt{#2}}% command name
    }%
    {% add the command and package to the index
      \index{#2 command@\protect\hangleft{\texttt{\tuftebs}}\texttt{#2} (\texttt{#1} package)}% command name
      \index{#1 package@\texttt{#1} package}\index{packages!#1@\texttt{#1}}% package name
    }%
}% command name -- adds backslash automatically
\newcommand{\docopt}[1]{\ensuremath{\langle}\textrm{\textit{#1}}\ensuremath{\rangle}}% optional command argument
\newcommand{\docarg}[1]{\textrm{\textit{#1}}}% (required) command argument
\newenvironment{docspec}{\begin{quotation}\ttfamily\parskip0pt\parindent0pt\ignorespaces}{\end{quotation}}% command specification environment
\newcommand{\docenv}[1]{\texttt{#1}\index{#1 environment@\texttt{#1} environment}\index{environments!#1@\texttt{#1}}}% environment name
\newcommand{\docenvdef}[1]{\hlred{\texttt{#1}}\label{env:#1}\index{#1 environment@\texttt{#1} environment}\index{environments!#1@\texttt{#1}}}% environment name
\newcommand{\docpkg}[1]{\texttt{#1}\index{#1 package@\texttt{#1} package}\index{packages!#1@\texttt{#1}}}% package name
\newcommand{\doccls}[1]{\texttt{#1}}% document class name
\newcommand{\docclsopt}[1]{\texttt{#1}\index{#1 class option@\texttt{#1} class option}\index{class options!#1@\texttt{#1}}}% document class option name
\newcommand{\docclsoptdef}[1]{\hlred{\texttt{#1}}\label{clsopt:#1}\index{#1 class option@\texttt{#1} class option}\index{class options!#1@\texttt{#1}}}% document class option name defined
\newcommand{\docmsg}[2]{\bigskip\begin{fullwidth}\noindent\ttfamily#1\end{fullwidth}\medskip\par\noindent#2}
\newcommand{\docfilehook}[2]{\texttt{#1}\index{file hooks!#2}\index{#1@\texttt{#1}}}
\newcommand{\doccounter}[1]{\texttt{#1}\index{#1 counter@\texttt{#1} counter}}

% Math
\usepackage{amsmath}  % extended mathematics
\usepackage{amsfonts}  % extended mathematics
\usepackage{amssymb}  % extended mathematics
\usepackage{amsthm}

\newtheorem{theorem}{Theorem}
\newtheorem{corollary}{Corollary}
\newtheorem{example}{Example}
\newtheorem{problem}{Problem}
\newtheorem{proposition}{Proposition}
\newtheorem{lemma}[theorem]{Lemma}
\newtheorem{definition}{Definition}

%%%%%%%%%%%%%%%%%%%%%%%%%%%%%%%%%%%%%%%%%%%%%%%%%%%%%%%%%%%%%%%%%%%%%%%%%%%%%%%
% Snippets

\newcommand{\snippetquadraticdicerule}{%
    For any given inequality $ax^2 + bx + c <> 0$, having
    $x_1$, $x_2$ (with $x_1 \leq x_2$) be
    the points such that $ax^2 + bx + c = 0$; the solution is
    $x < x_1 \wedge x > x_2$ if $a > 0$ and the sign is $>$ or if $a < 0$ and the sign is $<$;
    while the solution is $x_1 < x < x_2$
    if $a > 0$ and the sign is $<$ or if $a < 0$ and the sign is $>$.
}

\newcommand{\snippetdelopitaltheorem}{%
    When two differetiable
    functions $f$ and $g$ are
    such that:
    \begin{equation*}
        \lim_{x \to c} \frac{f(x)}{g(x)} = \left[\frac{\infty}{\infty}\right]
    \end{equation*}
    Then the following holds:
    \begin{equation*}
        \lim_{x \to c} \frac{f(x)}{g(x)} = \lim_{x \to c} \frac{f^\prime(x)}{g\prime(x)}
    \end{equation*}
    Note that $c$ can either be a number or $\infty$.
}


%%%%%%%%%%%%%%%%%%%%%%%%%%%%%%%%%%%%%%%%%%%%%%%%%%%%%%%%%%%%%%%%%%%%%%%%%%%%%%%

% Generates the index
\usepackage{makeidx}
\makeindex

%%%%%%%%%%%%%%%%%%%%%%%%%%%%%%%%%%%%%%%%%%%%%%%%%%%%%%%%%%%%%%%%%%%%%%%%%%%%%%%

\begin{document}

% Front matter
\frontmatter

% Blank page before title
%\blankpage

% Full title page
\maketitle

% Copyright page
\newpage
\begin{fullwidth}
~\vfill
\thispagestyle{empty}
\setlength{\parindent}{0pt}
\setlength{\parskip}{\baselineskip}
The 2nd Trial by \thanklessauthor

\par\smallcaps{Series: \thanklesspublisher}

\par\smallcaps{github.com/andry-tino/the-second-trial}

\par The 2nd Trial by Andrea Tino is licensed under
Attribution-NonCommercial-ShareAlike 4.0 International.
To view a copy of this license, visit
\url{http://creativecommons.org/licenses/by-nc-sa/4.0/}.\index{license}

\par\textit{First printing, \monthyear}
\end{fullwidth}

%%%%%%%%%%%%%%%%%%%%%%%%%%%%%%%%%%%%%%%%%%%%%%%%%%%%%%%%%%%%%%%%%%%%%%%%%%%%%%%

% Contents
\tableofcontents
%\listoffigures
%\listoftables

% Introduction
\cleardoublepage
%%%%%%%%%%%%%%%%%%%%%%%%%%%%%%%%%%%%%%%%%%%%%%%%%%%%%%%%%%%%%%%%%%%%%%%%%%%%%%%
% The 2nd Trial by Andrea Tino is licensed under CC BY-NC-SA 4.0. To view a   %
% copy of this license, visit                                                 %
% http://creativecommons.org/licenses/by-nc-sa/4.0/                           %
%%%%%%%%%%%%%%%%%%%%%%%%%%%%%%%%%%%%%%%%%%%%%%%%%%%%%%%%%%%%%%%%%%%%%%%%%%%%%%%

\chapter*{Introduction}

This sample book discusses the design of Edward Tufte's
books\cite{Tufte2001,Tufte1990,Tufte1997,Tufte2006}
and the use of the \doccls{tufte-book} and \doccls{tufte-handout} document classes.


%%%%%%%%%%%%%%%%%%%%%%%%%%%%%%%%%%%%%%%%%%%%%%%%%%%%%%%%%%%%%%%%%%%%%%%%%%%%%%%

\mainmatter

%%%%%%%%%%%%%%%%%%%%%%%%%%%%%%%%%%%%%%%%%%%%%%%%%%%%%%%%%%%%%%%%%%%%%%%%%%%%%%%
% The 2nd Trial by Andrea Tino is licensed under CC BY-NC-SA 4.0. To view a   %
% copy of this license, visit                                                 %
% http://creativecommons.org/licenses/by-nc-sa/4.0/                           %
%%%%%%%%%%%%%%%%%%%%%%%%%%%%%%%%%%%%%%%%%%%%%%%%%%%%%%%%%%%%%%%%%%%%%%%%%%%%%%%

\chapter[Study of functions]{Study of two functions}
\label{ch:fun}
\marginnote{Appears as \textit{first problem} in:
LI02, EA02 – Scientific; LI03 - Scientific - Applied Sciences;
LI15 - Scientific - Sport Sciences (standard session)}

This year as well we can find the classic study of function problem. We
are given 2 functions:
\begin{equation*}
    \begin{array}{cc}
        f(x) = ax^2 - x + b & g(x) = (ax + b)e^{2x - x^2} \\
    \end{array}
\end{equation*}
Before moving straight on with the questions the problem poses, let's just have
a look at $f$ and $g$ and
classify\sidenote{You need to understand who your enemy is if you
want to fight him and win!}
them. $f$ is a simple algebraic function: a second-order polynomial;
while $g$ is a more sophisticated thing: a transcendental function made up
by the combination of a first-order polynomial and an exponential term. Both
are parametric functions where, in addition to the indipendent variable
$x$, we can see parameters $a$ and $b$.

\paragraph[Problem 1]{Max and min}
\label{par:subp1_1}
\marginnote{\smallcaps{\nameref{par:subp1_1}}}
Considering $a \in \mathbb{R} \setminus \{ 0 \}$ and $b \in \mathbb{R}$, we
need to show that $g$ has a minimum and maximum (both absolute). I know at
this point you might be tempted to start calculating the first-order
derivative and the studying its sign, but you need to hold that thought
because we can solve this much more
quickly\sidenote{This is not just quicker though, it is much smarter: remember that
the problem does not ask to compute those extema, but to just prove their
existence.} without getting our hands too dirty. 

This is a maximization/minimization problem, therefore we need to
dig in our memory and find out what theorems can help us on this matter.
Weierstrass comes to mind: if we prove that $g$ is continuous inside a
closed interval, then $g$ will attain one minimum and one maximum;
however we do not have a closed interval here, we have the whole
$]-\infty, +\infty[$ so that is not a viable solution.
But we are on the right track here: \textbf{continuity is the key}.
If we can prove that $g$ is \smallcaps{continuous} and
\smallcaps{bounded}\sidenote{A function
$f(x) : \mathbb{R} \mapsto \mathbb{R}$ is bounded in
$x \in [a, b] \subset \mathbb{R}$ when there exist $m, M \in \mathbb{R}$
such that: $m \leq f(x) < M$. Note that the interval can also be
$x \in ]-\infty, +\infty[$.}
in $\mathbb{R}$, then we have
proved\sidenote{A more rigorous proof to why continuous bounded
functions attain extrema can be found at: \nameref{sec:contbound}.}
it has
absolute extrema. Think about it: continuity means that you can draw
$g$ without ever lifting your pen from the paper, the only thing
which can prevent the existence of extrema is if
$\lim_{x \to \infty} g(x) = \infty$.

Let's implement this strategy. Let's first check the domain of $g$
to see whether it is continuous. As stated earlier, this is the product
between two factors: a polynomial (always continuous in $\mathbb{R}$),
and an exponential (also continuous in $\mathbb{R}$). The argument
inside the exponential is another polynomial as well. Therefore
$g$ is defined in all $\mathbb{R}$ without any discontinuities.
We now need to prove $g$ is bounded, let's start on the left:
\begin{equation*}
    \lim_{x \to -\infty} g(x) = \lim_{x \to -\infty} (ax + b)e^{2x - x^2}
\end{equation*}
Remember that $a \neq 0$ which means tht the $x$ inside $ax + b$ is
always gonna be there, hence $\lim_{x \to -\infty} (ax + b) = -\infty$.
And of course the exponential will vanish as its argument tends
to $-\infty$: $\lim_{x \to -\infty} e^{2x - x^2} = 0$. We reach
an undecided form:
\begin{equation*}
    \lim_{x \to -\infty} (ax + b)e^{2x - x^2} =
        [-\infty \cdot 0]
\end{equation*}
Solving this indeterminate form can be done quickly by turning it into
a different one. Remember that every product $\alpha \beta$ can be
rewritten as $\frac{\alpha}{\beta^{-1}}$. What's the gain you ask?
If $\alpha \to \infty$ and $\beta \to 0$, then we have transformed a
$[\infty \cdot 0]$ form into a $[\frac{\infty}{\infty}]$.
\begin{equation*}
    \lim_{x \to -\infty} (ax + b)e^{2x - x^2} =
    \lim_{x \to -\infty} \frac{ax + b}{e^{x^2 - 2x}} =
    \left[\frac{\infty}{\infty}\right]
\end{equation*}
What good is this for? Well, now we can use De L\^{o}pital's
theorem\sidenote{When two differetiable
functions $f$ and $g$ are
such that:
\begin{equation*}
    \lim_{x \to c} \frac{f(x)}{g(x)} = \left[\frac{\infty}{\infty}\right]
\end{equation*}
Then the following holds:
\begin{equation*}
    \lim_{x \to c} \frac{f(x)}{g(x)} = \lim_{x \to c} \frac{f^\prime(x)}{g\prime(x)}
\end{equation*}
Note that $c$ can either be a number or $\infty$.}:
\begin{equation*}
    \lim_{x \to -\infty} \frac{ax + b}{e^{x^2 - 2x}} =
    \lim_{x \to -\infty} \frac{(ax + b)^\prime}{\left( e^{x^2 - 2x} \right)^\prime} =
    \lim_{x \to -\infty} \frac{a}{(2x - 2)e^{x^2 - 2x}} =
    \left[\frac{1}{\infty}\right] = 0
\end{equation*}
When we move on to take the limit of $g$ to $+\infty$,
the calculations are basically
the same and so is the result: $\lim_{x \to +\infty} g(x) = 0$.
That means the $g$ attains a maximum and a minimum because we just proved this
function is continuous and bounded.

\paragraph[Problem 2]{Parametric problem}
\label{par:subp2_1}
\marginnote{\smallcaps{\nameref{par:subp2_1}}}
We now need to find the values of $a$ and $b$ such that
$f$ and $g$'s graphs intersect at point $A \equiv (A_x, A_y) \equiv (2,1)$.
This is nothing new, we know that by putting $f$ and $g$ in a system, we can
find the values of $x$ such that the two functions will have the same value:
\begin{equation*}
    \begin{cases}
        f(A_x) = A_y\\
        g(A_x) = A_y
    \end{cases}
    \implies
    \begin{cases}
        aA_x^2 - A_x + b = A_y\\
        (aA_x + b)e^{2A_x - A_x^2} = A_y
    \end{cases}
    \implies
    \begin{cases}
        4a + b = 3\\
        2a + b = 1
    \end{cases}
\end{equation*}
By subtracting the first equation to the second, we get $2a = 2$, so:
$a = 1$ and, by substitution back in either two of the equations: $b = -1$.

\paragraph[Problem 3]{Drawing functions}
\label{par:subp3_1}
\marginnote{\smallcaps{\nameref{par:subp3_1}}}
Assuming $a = 1$ and $b = -1$, we want to draw $f$ and $g$ which become:
\begin{equation*}
    \begin{array}{cc}
        f(x) = x^2 - x -1 & g(x) = (x - 1)e^{2x - x^2} \\
    \end{array}
\end{equation*}
Let's start from $f$ which is very easy: it's a polynomial. Polynomials
are great becase we already know so many things upfront: they are
fully defined\sidenote{A function's \textit{definition} is another way
of referring to its domain.} in $\mathbb{R}$, they are not
bounded\sidenote{A function's \textit{boundedness} is a property related to
its codomain. When a function is unbounded, it means its codomain has no
restrictions, therefore the function takes values on the whole $y$ axis.},
and they diverge at $\infty$. This means
that: $f : \mathbb{R} \mapsto \mathbb{R}$, and that:
$\lim_{x \to \infty} f = \infty$. Plus, we can see $f$ is of second order,
which means it is a parabola with its concavity pointing
upwards\sidenote{A parabola $ax^2 + bx + c$ points upwards
when $a > 0$, downwards if $a < 0$.}. To draw a parabola, we only need
to find out its vertex and its intersection points with the axes.
To find the vertex you can either remember (from Geometry) its coordinates,
or use pure Calculus and treat this as a
minimization\sidenote[][0.05\textwidth]{Because the vertex
of an upward parabola happens
to also be the minimum.} problem (which means calculating the derivative
and study its sign). We are going for the second approach, but we will use
the generic equation of a parabola: $p(x) = ax^2 + bx + c$ to
find the generic coordinates of its vertex. We need to compute the first
order derivative: $p^\prime(x) = 2ax + b$ and equate it to $0$ to find
the extrema: $p^\prime(x) = 0$, which means: $2ax + b = 0$, which yields:
$x = -\frac{b}{2a}$. There is no need to study the sign of the derivative
because this is a parabola, and it does not have an horizontal inflection
point, so the $x$ we found is for sure an extremum.
To get the y-coordinate of the vertex, we simply
compute\marginnote[0.17\textwidth]{Do not forget that $\Delta = b^2 - 4ac$.}:
\begin{equation*}
    \begin{split}
        p\left( -\frac{b}{2a} \right) &=
        a \left( -\frac{b}{2a} \right)^2 + b \left( -\frac{b}{2a} \right) + c =
        a \frac{b^2}{4a^2} - \frac{b^2}{2a} + c\\
        &= \frac{b^2}{4a} - \frac{b^2}{2a} + c = \frac{b^2 - 2b^2 + 4ac}{4a} =
        - \frac{b^2 - 4ac}{4a} = - \frac{\Delta}{4a}
    \end{split}
\end{equation*}
So, even if you did not remember
the coordinates of the vertex
of a parabola, now you have
it: $V \equiv \left( -\frac{b}{2a}, - \frac{\Delta}{4a} \right)$.
After replacing\sidenote{From the definition of a
generic plynomial: $p(x) = ax^2 + bx + c$, and since
$f(x) = x^2 - x - 1$, we get that $a=1$, $b=-1$ and $c=-1$.}
$a$, $b$ and $c$ from the definition of $f$, we get its vertex:
$V_f \equiv \left( \frac{1}{2}, - \frac{5}{4} \right)$.
Let's now find where $f$ intersects the axes so we can get a total of 3 points
which is enough to draw the parabola. To find where $f$ crosses the x-axis,
we must put $y = f(x)$ in a system with $y = 0$ (x-axis' graph), which leads to
equation\marginnote[0.1\textwidth]{Note: $\sqrt{5} \approx 2.2$.}:
\begin{equation*}
    x^2 - x - 1 = 0 \implies x = \frac{-b \pm \sqrt{\Delta}}{2a} =
    \frac{1 \pm \sqrt{5}}{2} \implies
    \begin{cases}
        x_1 \approx 1.6\\
        x_2 \approx -0.6
    \end{cases}
\end{equation*}

\begin{figure*}[h] % Full width figure
    \begin{tikzpicture}
    \fill[Gray] (0,0) rectangle (16.7,4);
\end{tikzpicture}%
    \caption{This is the caption and much more to add come in here.}%
    \label{fig:fplot}%
\end{figure*}

So now we know $f$ also passes through points $F_1 \equiv (-0.6,0)$ and
$F_2 \equiv (1.6,0)$ other than $V_f$. It was expected to find two points
where $f$ crosses the x-axis because the vertex $V_f$ is negative, therefore
in the lower region of the Cartesian Plot. Note that we do not really need
to find where the function crosses the y-axis, 3 points is
enough\sidenote{Also remember that given 3 different points, there is only
one parabola crossing them all. That is why 3 points are sufficient for drawing.}.

\paragraph{}
Let's move to function $g$, this is gonna take a little more time. As
always, let's start from the domain: since $g$ is the product between
a polynomial (always defined) and an exponential (also always defined),
and considering that the argument of the exponential is another polynomial
(same story: always defined), we can clearly see that
$g : \mathbb{R} \mapsto \mathbb{R}$; there are no discontinuities.
Let's move on and try to find any
symmetries\sidenote{A function is \textit{even} when symmetric with
respect to the y-axis (the function is mirrored). While a function
is \textit{odd} when it is symmetric with respect to the plot's
center $O \equiv (0,0)$.}.
Why evaluating symmetry? Because it makes drawing a function quicker.
If we find out that a function is symmetric,
we just focus on half of the Cartesian Plot, the other half can be
filled using symmetry rules.
We must first compute $g(-x)$:
\begin{equation}\label{eq:gminusx}
    g(-x) = (-x - 1)e^{-2x - (-x)^2} = -(x + 1)e^{-2x - x^2}
\end{equation}
If $g$ is even, then $g(-x) = g(x)$. But we can see from equation
\ref{eq:gminusx} that is not the case.
If $g$ is odd, then $-g(-x) = g(x)$; let us compute $-g(-x)$:
\begin{equation*}
    -g(-x) = -(-x - 1)e^{-2x - (-x)^2} = (x + 1)e^{-2x - x^2}
        = (x + 1)e^{-(2x + x^2)}
\end{equation*}
Unfortunately $g$ is not odd either. Too bad, bad at least our effort was
still not in vain: we know something more about $g$ anyway.
Let's continue getting more info on $g$: let's study its sign:
\begin{equation}\label{eq:ggeqzero}
    g(x) \geq 0 \implies (x - 1)e^{2x - x^2} \geq 0
\end{equation}
Let's focus on the points where $g(x) = 0$. From equation
\ref{eq:ggeqzero} we see that the function is a product, and a
product is null when either two of its factors is null. But one of the
factors is an exponential which is never $0$, so we are left
with the binomial $x-1$ which becomes null at $x = 1$. Now we know
that $g$ crosses the x-axis in $x = 1$, which means it changes sign
there, because we proved that $g$ is continuous. Back
to equation \ref{eq:ggeqzero}, we can see that one of the factors
is an exponential which is always (strictly) positive,
therefore the sign of $g$
is decided by $x-1$ which is positive when $x > 1$:
\begin{equation*}
    g(x) > 0 \iff \forall x > 1
\end{equation*}
Moving on, we want to study $g$'s monotonicity. However, before digging
into calculations, remember that, in \nameref{par:subp1_1}, we have understood
$g$ is bounded and has at least one maximum and one minimum. Let's find them. 
We evaluate the first-order derivative:
\begin{equation}\label{eq:gprimedef}
    \begin{split}
        g^\prime(x) &= \frac{d}{dx} (x - 1) \cdot e^{2x - x^2} + (x - 1) \cdot \frac{d}{dx} e^{2x - x^2}\\
            &= e^{2x - x^2} + (x - 1)  e^{2x - x^2} \cdot \frac{d}{dx} (2x - x^2)\\
            &= e^{2x - x^2} + (x - 1) (2 - 2x) e^{2x - x^2}\\
            &= \left[ 1 + (x - 1) (2 - 2x) \right] e^{2x - x^2} = (-2x^2 + 4x - 1) e^{2x - x^2}
    \end{split}
\end{equation}
We can see $g^\prime$ is also the product between a polynomial and an exponential.
As we study its sign:
\begin{equation}\label{eq:gprimegeqzero}
    g^\prime(x) \geq 0 \iff (-2x^2 + 4x - 1) e^{2x - x^2} \geq 0
\end{equation}
We first want to spot were possible maxima and minima are located: $g^\prime(x) = 0$.
But, again, an exponential is never $0$, therefore we just need to focus
on\marginnote[0.1\textwidth]{Note: $\sqrt{2} \approx 1.4$.}:
\begin{equation}\label{eq:grpimenull}
    \begin{split}
        -2x^2 + 4x - 1 = 0 &\implies x = \frac{-b \pm \sqrt{\Delta}}{2a} =
        \frac{-4 \pm \sqrt{8}}{-4} = \frac{4 \pm 2\sqrt{2}}{4}\\
        &\implies
        x = 1 \pm \frac{\sqrt{2}}{2} \implies
        \begin{cases}
            x_1 = 1 + \frac{\sqrt{2}}{2} \approx 1.7\\
            x_2 = 1 - \frac{\sqrt{2}}{2} \approx 0.3
        \end{cases}
    \end{split}
\end{equation}
At this point we know that $g$ has possible\sidenote{Remember that when you find the
values for which a function's derivative is $0$, you have not found extrema:
you have found \smallcaps{candidate extrema}. You must study the derivative's sign
in order to confirm whether those candidates are maxima, minima or horizontal
inflections.} extrema in $x_1$ and $x_2$. Back to equation
\ref{eq:gprimegeqzero}, let's study $g^\prime$'s sign.
Since the exponential term is alway (strictly) positive, the overall sign is dictated
by $-2x^2 + 4x - 1$. But, in equation \ref{eq:grpimenull}, we already studied
this polynomial's zeroes\sidenote{Remember that when solving a
second-order algebraic inequality $ax^2 + bx + c <> 0$,
you must first evaluate the underlying equation: $ax^2 + bx + c = 0$.} in
equation \ref{eq:gprimegeqzero}, so we just apply the well known
rule\sidenote{For any given inequality $ax^2 + bx + c <> 0$, having
$x_1$, $x_2$ (with $x_1 \leq x_2$) be
the points such that $ax^2 + bx + c = 0$; the solution is
$x < x_1 \wedge x > x_2$ if $a > 0$ and the sign is $>$ or if $a < 0$ and the sign is $<$;
while the solution is $x_1 < x < x_2$
if $a > 0$ and the sign is $<$ or if $a < 0$ and the sign is $>$.}
for second-order algebraic inequalities and find out that: $g^\prime(x) > 0$ when
$-2x^2 + 4x - 1 > 0$ which holds when:
\begin{equation}\label{eq:gprimesignres}
g^\prime(x) > 0 \iff \forall x: x_1 < x < x_2
\end{equation}
This means that $g$ is monotonic incresing in between $x_1$ and $x_2$, decreasing
elsewhere. This makes sense because $g$ is bounded as we discovered in
\nameref{par:subp1_1}, which means it must go to $0$ when $x \to \pm\infty$:
$g$ decreases while we move further away towards $\infty$.
Becuase of equation \ref{eq:gprimesignres}, we can see that in $x_1$ and $x_2$ the sign
of $g^\prime$ gets inverted, which means that those two points are
extrema\sidenote{Whose existence we already proved earlier.}; let's calculate those:
\begin{equation*}
    \begin{cases}
        g(x_1) = \frac{\sqrt{2}}{2} e^{2\left( 1 + \frac{\sqrt{2}}{2} \right) - \left( 1 + \frac{\sqrt{2}}{2} \right)^2}
            = \sqrt{\frac{e}{2}} \approx 1.2\\
        g(x_2) = - \frac{\sqrt{2}}{2} e^{2\left( 1 - \frac{\sqrt{2}}{2} \right) - \left( 1 - \frac{\sqrt{2}}{2} \right)^2}
            = - \sqrt{\frac{e}{2}} \approx -1.2
    \end{cases}
\end{equation*}
And we have found the maximum: $M_g \equiv \left( 1 + \frac{\sqrt{2}}{2}, \sqrt{\frac{e}{2}} \right)$ and
the minimum: $m_g \equiv \left( 1 - \frac{\sqrt{2}}{2}, -\sqrt{\frac{e}{2}} \right)$.

At this point, normally, we would move on studying
the concavity by calcualting the second-order
derivative $g^{\prime\prime}$ and studying its sign;
but we can be smart and avoid all that calculation.
The key to this shortcut is understanding that $g$ is a
\smallcaps{smooth}\sidenote{A function $f$ is said to
be $n$-\textit{smooth}, and is written: $f \in \mathcal{C}^n$, when its $n$-th derivative is continuous. Now recall that
if the $n$-th derivative exists, so must the $(n-1)$-th and so forth. Also remember that if a function is differentiable,
then it is also continuous. It means that a $\mathcal{C}^n$ function has all its
first $n$ derivatives and they are continuous.} function. Smooth functions are special functions that present
no irregularities in their plot. If $g$ is at least 2-smooth, then its plot would be regular without
criticalities, spikes or wedges.
That means $g$ would start from $-\infty$ and decrease towards $m_g$ by changing
concavity before reaching $x_1$; the same goes when we start from $+\infty$ and move left towards $x_2$:
$g$ would change concavity before reaching $M_g$. Between $x_1$ and $x_2$ another concavity change would be
necessary because $g$ would not allow wedges. This approach sounds quick, so let's try to
prove $g$ is 2-smooth. We have already proved that $g$ is continuous; if we look at
equation \ref{eq:gprimedef}, $g^\prime$ is always the product between a polynomial and an exponetial:
both continuous functions, so $g^\prime$ is also continuous.
What about $g^{\prime\prime}$? We don't need to compute it to see it is continuous.
Look at $g^\prime$, differentiating it is the same process of
differentiating $g$ in equation \ref{eq:gprimedef}: you always end up with polynomials and exponential.
Therefore $g^{\prime\prime}$ is also
continuous\sidenote[][-0.2\textwidth]{Fun fact: if you extend
this line of thinking, you can prove that $g$ is $\mathcal{C}^2$ and above. Since every
further derivative will always result is more polynomials and exponentials, basically all $g$'s derivatives
are continuous. That is why $g \in \mathcal{C}^\infty$. Such functions are
called \smallcaps{infinitely differentiable}.}.

\begin{figure*}[h] % Full width figure
    \begin{tikzpicture}
    \fill[Gray] (0,0) rectangle (16.7,4);
\end{tikzpicture}%
    \caption{This is the caption and much more to add come in here.}%
    \label{fig:fplot}%
\end{figure*}

\paragraph[Problem 4]{Symmetry of $g$}
\label{par:subp4_1}
\marginnote{\smallcaps{\nameref{par:subp4_1}}}
By looking at $g$'s plot, we can assess the function had indeed a centre of symmetry
in $S_g \equiv (1,0)$. This evidence alone is quite sufficient to move on, but if we really wanna
solve the problem as asked, we have to prove $S_g$ is a centre of symmetry to $g$.
If a function $\phi(x)$\sidenote{We cannot use symbol $f$ because it's already taken.}
is symmetric with respect to a generic point whose x-coordinate is $x_S$, then the
following holds:
\begin{equation}\label{eq:fcsymm}
    \phi(x_S + x) = - \phi(x_S - x)
\end{equation}
So let's impose $x_S = 1$ and rewrite
equation \ref{eq:fcsymm}\sidenote{This equation and the technique
to find a generic centre of symmetry is
described at: \nameref{sec:fsymmcentre}.} as $g(1 + x) = - g(1 - x)$ and proceed:
\begin{equation*}
    \begin{split}
        (1 + x - 1)e^{2(1 + x) - (1 + x)^2} &= - (1 - x - 1)e^{2(1 - x) - (1 - x)^2} \implies\\
        xe^{2 + 2x - (1 + x^2 + 2x)} &= xe^{2 - 2x - (1 + x^2 - 2x)} \implies\\
        xe^{2 + 2x - 1 - x^2 - 2x} &= xe^{2 - 2x - 1 - x^2 + 2x)} \implies\\
        xe^{1 - x^2} &= xe^{1 - x^2)}
    \end{split}
\end{equation*}

\paragraph[Problem 5]{Tangency between $f$ ad $g$}
\label{par:subp5_1}
\marginnote{\smallcaps{\nameref{par:subp5_1}}}
From the plot of $f$ and $g$ we should see the two functions are tangent. We need to prove
the tangent point exists and it is $B \equiv (B_x,B_y) \equiv (0,-1)$.
Here's the strategy:
\begin{enumerate}
    \item \label{it:binfg} We substitute $B$ in both $f$ and $g$ verifying it belongs to both of them.
    \item \label{it:binderfg} To prove that $f$ and $g$ are tangent in $B$, we calculate
        both functions' derivatives in $B$ and assess they are the
        same\sidenote{Remember that the derivtive of a function represents
        the angular coefficient of the line tangent the function in each of its points.
        If two functions are tangent in a point, their derivatives in that point
        must be the same.}:
        $f^\prime(B_x) = g^\prime(B_x)$.
\end{enumerate}
Focusing on point \ref{it:binfg}, let's verify $B$ belongs to $f$:
\begin{equation*}
    B_y = f(B_x) \implies -1 = f(0) \implies (0)^2 - 0 - 1 \implies -1 = -1
\end{equation*}
That's ok. Let's check $B$ also belongs to $g$:
\begin{equation*}
    B_y = g(B_x) \implies -1 = g(0) \implies (-1)e^{0} \implies -1 = -1
\end{equation*}
Moving to point \ref{it:binderfg}, there's not much to do as we have already computed
$g^\prime$ and $f^\prime$ is very easy to calculate:
\begin{equation*}
    f^\prime(B_x) = g^\prime(B_x) \implies f^\prime(0) = g^\prime(0)
        \implies -1 = (-1) e^{0} \implies -1 = -1
\end{equation*}
And that proves $f$ and $g$ are tangent in $B$.

\newpar{Problem 6}{Area calculation}{par:subp6_1}
We are now tasked with finding the value of the area between the two curves $f$ and $g$.
We can clearly see from the plot we drew that they have two points in common;
since we already know one of them
is $B$, we need to find the other one which we will call $C = (C_x, C_y)$.
So we can use a system:
\begin{equation*}
    \begin{cases}
        y = f(x)\\
        y = g(x)
    \end{cases}
    \implies
    \begin{cases}
        y = x^2 - x - 1\\
        y = (x - 1)e^{2x - x^2}
    \end{cases}
\end{equation*}
Which leads to equation:
\begin{equation}\label{eq:tangeq}
    x^2 - x - 1 = (x - 1)e^{2x - x^2}
\end{equation}
Equation \ref{eq:tangeq} will be able to give us two solutions: one is $B_x$ and
the other one is $C_x$. But there is a problem: that is a
\smallcaps{polynomial-exponential}\sidenote{A
\textit{polynomial-exponential equation} is an equation where
$x$ or any of its powers appear both
inside and outside of one or more exponentials.} equation. Why it
is a problem? Because calculating the value of $x$ in
polynomial-exponential equations using symbolic calculus is
impossible\sidenote{You can find the proof to why polynomial-exponential
equations are impossible to solve analytically at: \nameref{sec:polyexpeqimp}.}!
Yeah, just forget it, you simply cannot do it; but then it means that the problem
of calculating $C$ is impossible as well? Of course not! If we cannot
find out $C$ analytically, we can try to do that \textbf{numerically}.
The numerical methods we know are all about finding zeroes of a function;
to get that function, we can start from equation \ref{eq:tangeq}: $f(x) = g(x)$,
and move all terms to the LHS: $f(x) - g(x) = 0$, and we can simply
introduce a new function $\varepsilon(x) = f(x) - g(x)$; the equation we get
is responsible for finding the zeroes of $\varepsilon$; one of those zeroes is
exactly what we are looking for: $C_x$.
The next step is choosing an interval $[x_1, x_2] \subset \mathbb{R}$: a range
of values of $x$ inside which we are going to iteratively try to find $C_x$.
We need to be sure that: $C_x \in [x_1, x_2]$.
By looking at the plot we drew before, we can see that $C_x$ comes after the point
where $f$ intersects the x-axis, so we can
set $x_1 = \frac{1 + \sqrt{5}}{2}$\sidenote{We calculated that point 
at: \nameref{par:subp3_1}.}. What about the other end of the interval?
Well that is tricky because we do not really have a reference on the plot to
guide us, so let's try to
eyeball it: we can pick up a number $x_2 > x_1$ and then calculate $f(x_2)$ and
$g(x_2)$, if $f(x_2) > g(x_2)$ then we know for sure we are after $C_x$ because
that is exactly where $f$ and $g$ meet and then $f$ starts becoming
greater than $g$. Let's try $x_2 = 3$: $f(3) = 3^2 - 3 - 1 = 5$, and
$g(3) = 2e^{-3} = \frac{2}{e^3}$, it is clear that $5 > \frac{2}{e^3}$, so we
can use $x_2 = 3$.

%%%%%%%%%%%%%%%%%%%%%%%%%%%%%%%%%%%%%%%%%%%%%%%%%%%%%%%%%%%%%%%%%%%%%%%%%%%%%%%
% The 2nd Trial by Andrea Tino is licensed under CC BY-NC-SA 4.0. To view a   %
% copy of this license, visit                                                 %
% http://creativecommons.org/licenses/by-nc-sa/4.0/                           %
%%%%%%%%%%%%%%%%%%%%%%%%%%%%%%%%%%%%%%%%%%%%%%%%%%%%%%%%%%%%%%%%%%%%%%%%%%%%%%%

\chapter[Magnetism and capacitors]{Magnetic field inside a
capacitor}
\label{ch:magn}
\marginnote{Appears as \textit{second problem} in:
LI02, EA02 – Scientific; LI03 - Scientific - Applied Sciences;
LI15 - Scientific - Sport Sciences (standard session)}

This problem introduces a flat capacitor to which
a specific voltage is applied. The voltage (which is null at $t = 0$)
varies over time and it causes a magnetic field $\vec{B}$ to originate,
whose module\sidenote{Disregarding border effects.} is:
\begin{equation}\label{eq:bdef}
    B = \frac{k t}{\sqrt{\left( t^2 + a^2 \right)^3}} \cdot r
\end{equation}
Having the following
quantities\sidenote[][0.17\textwidth]{The text does not mention what MU is
assigned to $r$, but it is sensible to
infer that it should be the same as $R$ because a constraint is present
between them: $r \leq R \implies r - R \leq 0$, and since they sum together,
they must have the same MU.}:

\begin{table}[h]
\begin{center}
    \footnotesize%
    \begin{tabular}{cccl}
    \toprule
    Quantity & UM & Constraints & Description \\
    \midrule
    $\vec{B}$ & \unit{T} & & Magnetic field \\
    $R$ & \unit{m} & $R > 0$ & Capacitors plate's radius \\
    $d$ & \unit{m} & $d > 0$ & Distance between plates \\
    $r$ & \unit{m} & $r > 0 \wedge r \leq R$ & Distance from axis \\
    $k$ & & $k > 0$ & Constant \\
    $a$ & & $a > 0$ & Constant \\
    $t$ & \unit{s} & $t \geq 0$ & Time \\
    \bottomrule
    \end{tabular}
\end{center}
\caption{Most important quantities in the problem.}
\label{tab:quants}
\end{table}

\begin{figure}
    \centering
    \begin{tikzpicture}
    \fill[Gray] (0,0) rectangle (10,3.5);
\end{tikzpicture}
    \caption[Problem representation.][0pt]{Problem representation: a 3D view and a 2D
    simplified diagram.}
    \label{fig:capacitorovervw}
\end{figure}

\paragraph[Problem 1]{What is the UM of $a$ and $k$?}
\label{par:subp1}
\marginnote{\smallcaps{\nameref{par:subp1}}}
We start off with a question
about \smallcaps{dimensionality}. By looking at equation \ref{eq:bdef}, we can
easily set up a dimensional equation:
\begin{equation*}
    \left[ B \right] = \left[ \frac{k t}{\sqrt{\left( t^2 + a^2 \right)^3}} r \right] \implies
    \unit{T} = \frac{[k] \cdot \unit{s} \cdot \unit{m}}{\sqrt{\left( \unit{s}^2 + [a]^2 \right)^3}}
\end{equation*}
Notice that, at denominator on the LHS, we have: $s^2 + [a]^2$. Since $[a]^2$
appears in a sum with $s^2$, it must have the same UM: $[a]^2 = s^2$.
That means $[a] = \unit{s}$, and the whole sum collapses into a single term: $s^2$:
\begin{equation*}
    \unit{T} = \frac{[k] \cdot \unit{s} \cdot \unit{m}}{\sqrt{\left( \unit{s}^2 \right)^3}} = 
    \frac{[k] \cdot \unit{s} \cdot \unit{m}}{\sqrt{\left( \unit{s}^3 \right)^2}} = 
    \frac{[k] \cdot \unit{s} \cdot \unit{m}}{\unit{s}^3} =
    \frac{[k] \cdot \unit{m}}{\unit{s}^2}
\end{equation*}
Notice how we further simplified the expression by removing the square root and removing
$\unit{s}$ terms. At this point, we have two choices: either we remember that
$T=\unit{Kg} \cdot \unit{s}^{-1} \cdot \unit{C}^{-1}$, or we try to calculate it from
Maxwell's
equations\sidenote{You can use Maxwell's 4th equation as described at: \nameref{sec:maxwell}.}.
Either way, we end up with:
\begin{equation*}
    \frac{\unit{Kg}}{\unit{s} \cdot \unit{C}} = \frac{[k] \cdot \unit{m}}{\unit{s}^2} \implies
    \frac{\unit{Kg}}{\unit{C}} = \frac{[k] \cdot \unit{m}}{\unit{s}} \implies
    \frac{\unit{Kg} \cdot \unit{s}}{\unit{C} \cdot \unit{m}} = [k]
\end{equation*}
Which means that $[k] = \unit{Kg} \cdot \unit{s} \cdot \unit{C}^{-1} \cdot \unit{m}^{-1}$.

\paragraph[Problem 2]{Why $B \neq 0$ in the capacitor even without magnets or currents?}
\label{par:subp2}
\marginnote{\smallcaps{\nameref{par:subp2}}}
The answer to this
is given by remembering that a magnetic field can originate from two possible sources:
\begin{itemize}
    \item Electrical currents.
    \item Variations of electrical fields.
\end{itemize}
If you remember Maxwell's 4th
equation\sidenote{See: \nameref{sec:maxwell}.},
then you can actually have a confirmation of these
two facts. In our problem, there are no magnets and there is no electrical current because a
capacitor opens the circuit blocking any current from flowing. However a voltage is applied,
that means an electrical field originates between the plates; and the voltage varies over
time which means the electrical field varies too, and that causes a magnetic field to
arise as per second point above.

\paragraph[Problem 3]{How do the directions of $\vec{E}$ and $\vec{B}$ relate together?}
\label{par:subp3}
\marginnote{\smallcaps{\nameref{par:subp3}}}
Here as well you should remember that inside an electromagnetic field, the electrical
field $\vec{E}$ and the magnetic field $\vec{B}$ are always perpendicular to each other.
If you didn't remember that, take Maxwell's either 3rd of 4th
equation\sidenote[][-0.1\textwidth]{The proof that $\vec{E}$ and $\vec{B}$ are perpendicular can be
found at: \nameref{sec:maxwell}.},
and you will be able to check it is true.

\paragraph[Problem 4]{Circuital current of $\vec{B}$}
\label{par:subp4}
\marginnote{\smallcaps{\nameref{par:subp4}}}
%
\begin{marginfigure}
    \centering
    \begin{tikzpicture}
    \fill[Gray] (0,0) rectangle (4,4);
\end{tikzpicture}
    \caption[Circuital current over a closed loop.]{Circuital current over a closed loop.}
    \label{fig:circuital}
\end{marginfigure}
%
We now want to calculate the circuital electrical current that $\vec{B}$ causes
over closed line $C$. This part is where we need to really remember the theory about
Electromagnetism. There is a precise sequence of events happening here, so let's
see what exactly happens between the capacitor's plates:
\begin{enumerate}
    \item A voltage is applied to the capacitor, causing electrical charges to accumulate on the
        plates.
    \item The presence of charges causes a uniform electrical field $\vec{E}$ to
        originate\marginnote{Maxwell's 1st law: electrical charges generate electrical fields.}.
    \item The voltage changes in time, that causes the amount of charge $Q$ on the plates to
        vary over time. This means that $\vec{E}$ changes as well over time.
    \item By virtue of Maxwell's 4th equation, we know that a changing electrical field
        causes a magnetic field $\vec{B}$ to arise.
    \item Ampere's law states that a current flows inside a closed
        line\marginnote{Also called: \textit{Amperian loop}.} permeated by a magnetic field.
        Loop $C$ is exactly that Amperian loop and $i_C$ is the so called
        \textit{displacement current} caused by $\vec{B}$.
\end{enumerate}
We want to
calculate\sidenote{At: \nameref{sec:maxwell}, I explain more in detail this process.} $i_C$!
The tool that can help us do that is Maxwell's 4th equation:
\begin{equation}
\label{eq:max4}
    \oint_C \vec{B} \cdot \,d \vec{l} = \mu_0 (i_{\text{cond}} + i_{\text{disp}})
\end{equation}
I know that is a scary equation, but I promise it's not that bad.
Equation \ref{eq:max4} states that
a magnetic field $\vec{B}$ can originate from two possible currents:
$i_{\text{cond}}$ represents the current passing through surface $S$ bounded by loop $C$;
while $i_{\text{disp}}$ is the current flowing in $C$ that originates from
the varying electrical field $\vec{E}$. So the current flowing in $C$ is the
displacement current, which means: $i_{\text{disp}} = i_C$.
So we can rewrite equation \ref{eq:max4} as:
\begin{equation}
    \label{eq:max4full}
    \oint_C \vec{B} \cdot \,d \vec{l} = \mu_0 (i_{\text{cond}} + i_C)
\end{equation}
We know that $i_C \neq 0$ because of $\vec{E}$ varying
in time. What about $i_{\text{cond}}$? That is the current that cuts through $S$. But
we are in the middle of the plates and there is no current flowing between the plates
which is flowing through $S$, that's why $i_{\text{cond}} = 0$. Thanks to these new
facts, we can now simplify equation \ref{eq:max4full} as:
\begin{equation}
    \label{eq:max4slim}
    i_C = \frac{1}{\mu_0}\oint_C \vec{B} \cdot \,d \vec{l}
\end{equation}
All right, now we really need to take care of the big elephant in the room we have
avoided looking at so far: the scary integral.
That integral is the sum of infinitely many terms: the scalar product between
the magnetic field $\vec{B}(x,y,z)$ at a certain point $(x,y,z)$ on the
loop $C$; and $\,d \vec{l}$ which the oriented piece of line $C$ at the same point.
By taking all points on $C$, calculating that scalar product for all of them, and
summing up all results, we get the value of the integral. But let's take a look at the
term $\vec{B} \cdot \,d \vec{l}$; remember that this quantity has to be evaluated for
every point on loop $C$. Being a scalar
product\sidenote{Remember that the scalar product between two vectors is a number:
    \begin{equation*}
        \vec{a} \cdot \vec{b} = a \cdot b \cdot cos\theta
    \end{equation*}
    Where $\theta$ is the angle between vectors $\vec{a}$ and $\vec{b}$.
    And it represents the length of the projection of $\vec{a}$ on $\vec{b}$.
}
we have that:
$\vec{B} \cdot \,d \vec{l} = B \,d l \cos\theta$.
What is the value of $\theta$? That is the angle between the oriented infinitesimal
loop line and the magnetic field vector. We know that the magnetic field is
perpendicular\sidenote{We discovered that at \nameref{par:subp3}.} to the
electrical field,


%%%%%%%%%%%%%%%%%%%%%%%%%%%%%%%%%%%%%%%%%%%%%%%%%%%%%%%%%%%%%%%%%%%%%%%%%%%%%%%
% The 2nd Trial by Andrea Tino is licensed under CC BY-NC-SA 4.0. To view a   %
% copy of this license, visit                                                 %
% http://creativecommons.org/licenses/by-nc-sa/4.0/                           %
%%%%%%%%%%%%%%%%%%%%%%%%%%%%%%%%%%%%%%%%%%%%%%%%%%%%%%%%%%%%%%%%%%%%%%%%%%%%%%%

\chapter[Questions]{Questions}
\label{ch:qs}

In this year's exams we can see a nice variaty of topics in the questions
ranging from the main theme: Magnetism to Probability and Linear Algebra.

\section{Extrema finding}
\label{sec:extremafind}
\marginnote{Appears as problem 1 in L 06.}
Let $p(x)$ be a polynomial, the following function is provided:
\begin{equation*}
    f(x) = \frac{p(x)}{x^2 + d}
\end{equation*}
Where $d \in \mathbb{R}$. $f$'s plot intersects the x-axis at $x_1 = 0$ and
$x_2 = \frac{12}{5}$, and has these lines as asymptotes: $x=3$, $x=-3$ and
$y=5$. We need to calculate the points of relative maximum and minimum.

\paragraph{}
The first thing important noticing is that $f$ is the ratio between two
polynomials. How much do we know about these two polynomials? The
denominator is a second-degree polynomial, while we do not have much
information about the numerator. We need to know a few things more about
$p$. From what the problem tells us, we know that
$\lim_{x \to \infty} f(x) \neq \infty$. That is a very important piece
of data because we know that the ratio between two polynomials tends
to a finite value (as the variable tends to infinity) only when they
have the same
degree\sidenote{The full theorem actually states that:
\begin{equation*}
    \lim_{x \to \infty} \frac{N(x)}{D(x)} = l \neq \infty
        \iff n = m
\end{equation*}
Where $N(x)$ and $D(x)$ are two polynomials with
degrees $n$ and $m$ respectively. And we also have that:
\begin{equation}
    l = \frac{a_N}{a_D}
\end{equation}
Where $a_N$ id $N$'s highest power's coefficient, and
$a_D$ id $D$'s highest power's coefficient.}.
Thanks to it, we can rewrite $f$ as:
\begin{equation}\label{eq:fracinfty}
    f(x) = \frac{ax^2 + bx + c}{x^2 + d}
\end{equation}
Because of equation \ref{eq:fracinfty}, we have that:
\begin{equation*}
    \lim_{x \to \infty} f(x) = \lim_{x \to \infty} \frac{ax^2 + bx + c}{x^2 + d}
        = \frac{a}{1} = 5
        \implies a = 5
\end{equation*}
As we unveil the values of coefficients, we know more and more the final form
of $f$, and this is important because we will need to compute its derivative and
then study its sign: the less parameters in the derivative, the easier our job.
Now $f$ is:
\begin{equation*}
    f(x) = \frac{5x^2 + bx + c}{x^2 + d}
\end{equation*}
We still have 3 parameters.
Let's use the info about intersections with the x-axis to find the values of some
of those:
\begin{equation*}
    \begin{split}
        &\begin{cases}
            f(0) = 0\\
            f\left( \frac{12}{5} \right) = 0
        \end{cases}
        \implies
        \begin{cases}
            \frac{c}{d} = 0\\
            \frac{5\frac{12^2}{5^2} + b\frac{12}{5} + c}{\frac{12^2}{5^2} + d} = 0
        \end{cases}
        \implies
        \begin{cases}
            c = 0\\
            \frac{\frac{12^2}{5} + b\frac{12}{5}}{\frac{12^2}{5^2} + d} = 0
        \end{cases}
        \implies\\
        &\frac{12^2 + 12b}{5} = 0 \implies
        12 + b = 0 \implies
        b = -12
    \end{split}
\end{equation*}
Which means now $f$ can be written as:
\begin{equation*}
    f(x) = \frac{5x^2 - 12x}{x^2 + d}
\end{equation*}
We only have $d$ left to find out, and for that we can use the information about
the two vertical asymptotes. Remember that $f$ is a ratio, when the denominator
tends to 0, the whole ratio tends to infinity; it means that if we find the
values of $x$ for which $x^2 + d = 0$, we will find the exact points where
the asymptots are: $x^2 + d = 0 \implies x = \pm \sqrt{-d}$. At this point
it becomes clear that $d < 0$ because otherwise we would have no asymptotes as
we would need to extract the square root of a negative number which is not a
real number.
Also notice that $d \neq 0$, otherwise we would not have two different
asymptotes, but only one. Since we know that the two aymptotes are $x = \pm 3$,
we can conclude that: $\pm 3 = \sqrt{-d} \implies d = -9$, and:
\begin{equation*}
    f(x) = \frac{5x^2 - 12x}{x^2 - 9}
\end{equation*}
Now we can compute $f^\prime$:
\begin{equation*}
    f^\prime(x) = \frac{10x - 12}{x^2 - 9} - \frac{5x^2 - 12x}{(x^2 - 9)^2} 2x = \dots =
        6\frac{2x^2 - 15x + 18}{(x^2 - 9)^2}
\end{equation*}
We first need to find the zeros\sidenote{Remember that the zeros
of a fraction $\frac{N(x)}{D(X)}$ can be found by finding the zeros of the numerator $N(x)$.}
of $f^\prime$:
\begin{equation*}
    \begin{split}
        2x^2 - 15x + 18 = 0 &\implies x = \frac{-b \pm \sqrt{\Delta}}{2a} =
        \frac{15 \pm \sqrt{81}}{4} = \frac{15 \pm 9}{4}\\
        &\implies
        \begin{cases}
            x_1 = \frac{15 + 9}{4} = \frac{24}{4} = 6\\
            x_2 = \frac{15 - 9}{4} = \frac{6}{4} = \frac{3}{2}
        \end{cases}
    \end{split}
\end{equation*}
Switching to the corresponding inequality: $2x^2 - 15x + 18 > 0$, we see that
the leading coefficient is positive and the inequality's sign is $>$; as
already stated in \ref{sn:dice}, we have
that: $f^\prime(x) > 0 \iff x < \frac{3}{2} \vee x > 6$.
%
\begin{marginfigure}[-5.3\textwidth]
    \includegraphics[width=1\linewidth]{qs_1_fplot_diag.tikz}
    \caption{Construction of the final plot.}%
    \label{fig:pr1fplot}%
\end{marginfigure}
%
This confirms that $x_1$ and $x_2$ are extrema (since around them $f^\prime$ changes sign),
and gives us the final result: $x_1$ is a point of relative maximum, while $x_2$ is
a point of relative minimum.

\section{A big polynomial}
\label{sec:bigpoly}
\marginnote{Appears as problem 1 in L 06.}
Let $g(x)$ be the following polynomial:
\begin{equation*}
    g(x) = \sum_{n = 1}^{1010} x^{2n-1} = x + x^3 + x^5 + \dot + x^{2017} + x^{2019}
\end{equation*}
We must prove that $g$ has a single real root. We also need to compute:
\begin{equation}\label{eq:limfracpolyexp}
    \lim_{x \to +\infty} \frac{g(x)}{1.1^x}
\end{equation}

\paragraph{}
The first part of the problem is about proving the \textbf{existence} and \textbf{uniqueness}
of $g$'s real root. In mathematical terms, this is what we need to prove:
\marginnote{You might see logical operator $\exists!$ for the first time: it
simply means: \textit{"there exists only one..."}.}
\begin{equation*}
    \exists! x_0 \in \mathbb{R} : g(x_0) = 0
\end{equation*}
The proof is actually much easier than you can expect. $g$ looks very scary because it's big,
but if we pay attention, we can see that the polynomial is missing the contant term $x^0$.
So we can actually rewrite $g$ by extracting $x$ as:
\begin{equation*}
    g(x) = x \cdot \left(1 + x^2 + x^4 + \dot + x^{2016} + x^{2018} \right)
\end{equation*}
Look now, we have basically applied the following transformation:
\begin{equation*}
    g(x) = \sum_{n = 1}^{1010} x^{2n-1} = x \cdot \sum_{n = 1}^{1009} x^{2n} 
\end{equation*}
At this point it's easy to see there exists only a single real root to $g$:
\begin{equation}\label{eq:gsumprodzero}
    g(x) = 0 \implies x \cdot \sum_{n = 1}^{1009} x^{2n} = 0
\end{equation}
And that is $x_0 = 0$:
\begin{enumerate}
    \item As equation \ref{eq:gsumprodzero} shows, $g$ is the product of two factors: $g$ is
        null when either one of them is null. One of these terms is expression $x$:
        therefore $x = 0$ is a real root. This proves the existence of $x_0$.
    \item The second factor in $g$'s definition as per equation \ref{eq:gsumprodzero},
        is a complete polynomial where all its coefficients are positive and all powers
        are even\sidenote{Expression $2n$ for $n = 1, 2, 3, \dots$ describes an even number.}.
        This means that the sum we see as second factor is a sum of squares, and
        we know that a sum of squares is never null (for real values). This proves the uniqueness
        of $x_0$.
\end{enumerate}

\paragraph{}
As for the second part of the problem, the limit in equation \ref{eq:limfracpolyexp} involves
a ratio between a polynomial and an exponential term. It is easy to prove that ratio tends
to $0$ when $x \to +\infty$:
\begin{enumerate}
    \item We know that an exponential term $a^x$ tends to $+\infty$ as the variable tends
        to $+\infty$ when $a > 1$. Therefore $1.1^x \to +\infty$.
    \item We know that polynomials always tend to $\infty$ when the variable tends to $\infty$.
        In this case, as $x \to +\infty$, $g(x) \to +\infty$ because the leading coefficient
        is positive.
    \item The ratio $\frac{g(x)}{1.1^x}$ tends to $0$ as $x \to +\infty$ because the exponential
        appears as the denominator. An exponential grows to $+\infty$ with an incredible speed.
        A polynomial cannot match that speed. Therefore the denominator grows much more than
        the numerator causing the whole fraction to tend towards $0$.
\end{enumerate}
If we want to be very precise and provide a proof to the last point, we can apply De L\^{o}pital's
theorem\textsuperscript{\ref{sn:dehopital}}. Both functions are
differentiable\sidenote{Actually they are both $\mathcal{C}^\infty$: indefinitely differentiable.},
and they both tend towards $\infty$; so we get the indeterminate
form $\left[ \frac{\infty}{\infty} \right]$. As we apply the rule once, we notice that the
indeterminate form still remains but the polynomial has dropped its order by 1; so we re-apply
the rule (as many time as the original order of the polynomial: $1010$) until we get $0$:
\marginnote{Each arrow indicates an iteration where we take the derivative of the numerator and
the denominator. Also note that $p^{[N]}(x)$ indicates a generic $N$-degree polynomial.}
\begin{equation*}
    \frac{p^{[N]}(x)}{1.1^x} \rightarrow \frac{p^{[N-1]}(x)}{\log(1.1)1.1^x}
        \rightarrow \frac{p^{[N-2]}(x)}{\log^2(1.1)1.1^x} \rightarrow \dots \rightarrow
        \frac{0}{\log^N(1.1)1.1^x} = 0
\end{equation*}
If we run the above process for $N=1010$, we get the result.

\section{Parallelepipeds}
\label{sec:paralpyd}
\marginnote{Appears as problem 1 in L 06.}
We must find, among all square-base parallelepipeds with total surface $S$, the one
which minimizes the sum of all sides.
%
\begin{marginfigure}
    \includegraphics[width=1\linewidth]{qs_3_parapip_diag.tikz}
    \caption{Characteristics of the parallelepiped.}%
    \label{fig:pr3parapip}%
\end{marginfigure}
%

\paragraph{}
This problem is very short but absolutely non-trivial. That doesn't mean it is
very difficult; however, requires a good setup at the beginning.
As every minimization scenario, we must formulate the problem into an
equation. Let's start by writing down the formulas for the total surface
and the sum of all sides:
\begin{equation*}
    \begin{array}{cc}
        S = 2l^2 + 4lh & L = 8l + 4h \\
    \end{array}
\end{equation*}
Here, $l$ is the length of one side of the square base, $h$ is the height of
the solid, and $L$ is the sum of all sides.
Now let's try to think a little. We want to impose a certain surface $S$,
it is clear that we can play with $l$ and $h$ to generate many different
parallelepipeds that will keep $S$ unchanged. In this scenario, if we, for example,
try to increase $h$, we will see the solid getting taller, but the base will
probably shrink ($l$ will decrease) because we want to keep $S$ constant.
So let's take the equation for $S$, there are 3 variables therein: $S$, $l$ and $h$.
Let's consider $S$ fixed from now on; it means that $S$ will be a known value that we
choose. So now we have only 2 variables: $l$ and $h$: let's calculate one as function of
the other. We choose (of course) to
write $h = f(l)$\sidenote{Because the equation is linear in
$h$ and we can quickly isolate that term. If we chose to isolate $l$ instead,
we would need to solve a quadratic equation which is more difficult.}:
\begin{equation}\label{eq:hasfofl}
    h = \frac{S}{4} \cdot \frac{1}{l} - \frac{l}{2}
\end{equation}
Equation \ref{eq:hasfofl} sets a contraint on $h$ based on the values of $l$. If
we respect equation \ref{eq:hasfofl}, the values of $h$ and $l$ will cause
the total surface of the solid to be $S$.
Let's now substitute $h$ from equation \ref{eq:hasfofl} into the equation of
$L$:
\begin{equation}\label{eq:Lasf}
    L(l) = 8l + 4h = 8l + 4\left( \frac{S}{4} \cdot \frac{1}{l} - \frac{l}{2} \right)
        = \dots = 6l + \frac{S}{l}
\end{equation}
In equation \ref{eq:Lasf} $L$ is purposely indicated as a function of $l$. Now that we
have the total sum depending on $l$, we can find the extrema by following the normal
minimization process: compute the derivative, find its zeros and study the sign.
\begin{equation}\label{eq:Lfder}
    L^\prime(l) = \frac{d}{dl} \left( 6l + \frac{S}{l} \right) = 6 - \frac{S}{l^2}
\end{equation}
Let's impose $L^\prime(l) = 0$:
\begin{equation}\label{eq:Lroots}
    6 - \frac{S}{l^2} = 0 \implies \frac{6l^2 - S}{l^2} = 0 \implies
        6l^2 - S = 0 \implies l = \pm \sqrt{\frac{S}{6}}
\end{equation}
Value $\frac{S}{6}$ is positive so we can actually extract its square
root\sidenote{We must always check that all values make sense, this means
that equation \ref{eq:Lroots} is legit.}. We get 2 values
for $l$: one positive and one negative.
Since $l$ must be positive (it is a length), we must discard the negative one and
only accept $l = +\sqrt{\frac{S}{6}}$. But we still must make sure this value is
an extremum and a point of minimum.
To check that, we look at equation \ref{eq:Lfder}: $6l^2 - S$ is a parabola
pointing upwards\sidenote{The leading coefficient is $6 > 0$.}. The root
we have taken is the greatest one: it is where the parabola crosses the x-axis
from the vertex to infinity; therefore we pass from negative values to positive values,
hence we found a point of minimum. Now that we have $l$, we can calculate the minimum sum of all
sides $L^\ast$ to have surface $S$, by substituting in equation \ref{eq:Lasf}:
\begin{equation*}
    L^\ast =  6 \left( \frac{S}{6} \right)^{\frac{1}{2}} + S \left( \frac{S}{6} \right)^{-\frac{1}{2}}
\end{equation*}
Note that we must add $L^\ast > 0$
as a legit condition.

\section{Spheres}
\label{sec:speheres}
\marginnote{Appears as problem 1 in L 06.}
We are asked to
prove that the geometric locus of points $P$ \textbf{in space} such that the ratio
of the distances between them and two arbitrary points $A$ and $B$ (with $A \neq B$)
is always $\sqrt{2}$, is actually a sphere. The problem poses $A \equiv (2,0,-1)$ and
$B \equiv (-2,2,1)$.

\paragraph{}
I have to admit that this problem really puzzled me, as I never heard of this alternative
definition of the sphere. So when I wrote down my calculations and actually found out it
is true, I was very surprised. Let's start with the classic
definition\sidenote{A sphere is the geometric locus of points $P$ such that they all
have the same distance from one special point called \textit{center}: $C$.} of
the sphere we all know:
\marginnote{We are indicating a generic point on the sphere as
$P \equiv (x,y,z)$, and the center of the sphere as
point: $C \equiv (x_0,y_0,z_0)$.}
\begin{equation*}
    d(P,C) = R \implies \sqrt{(x-x_0)^2 + (y-y_0)^2 + (z-z_0)^2} = R
\end{equation*}
If we raise to the second power\sidenote{Always remember that
when raising a square root to the second power you get the absolute value of the operand:
\begin{equation*}
    \sqrt{x^2} = \left( \sqrt{x} \right)^2 = |x|
\end{equation*}
If we can guarantee that the operand is always positive, then we can remove the
absolute value. That is the reason why we do it here: the quantity inside the
square root is a sum of squares.} both hands of the equation, we get:
\begin{equation}\label{eq:sphereq}
    (x-x_0)^2 + (y-y_0)^2 + (z-z_0)^2 = R^2
\end{equation}
Equation \ref{eq:sphereq} is the classic equation of a sphere.
Let's now have a look at the definition of sphere that the problem gives us.
If we define $A \equiv(x_1,y_1,z_1)$ and $B \equiv(x_2,y_2,z_2)$, we see that
such definition can be translates as: $\frac{d(P,A)}{d(P,B)} = \sqrt{2}$, which
becomes:
\begin{equation}\label{eq:sphereqalt}
    \frac{\sqrt{(x-x_1)^2 + (y-y_1)^2 + (z-z_1)^2}}{\sqrt{(x-x_2)^2 + (y-y_2)^2 + (z-z_0)^2}} = \sqrt{2}
\end{equation} 
Our strategy is simple: let's try to transform\sidenote{By \textit{transform} we mean
rearranging one equation to have the same \smallcaps{structure} of another.} equation \ref{eq:sphereqalt}
into equation \ref{eq:sphereq}; if we are successful, then we would solve the problem.
Let's start by raising both sides of the equation to the second power\sidenote{Here as well,
no need for the absolute value as the quantities under square root are all positive}:
\begin{equation}\label{eq:sphereqalt2}
    \frac{(x-x_1)^2 + (y-y_1)^2 + (z-z_1)^2}{(x-x_2)^2 + (y-y_2)^2 + (z-z_0)^2} = 2
\end{equation}
Let's have a look at equation \ref{eq:sphereq}: in there, each variable ($x$, $y$ and $z$)
is summed to a constant ($x_0$, $y_0$ and $z_0$); while in equation \ref{eq:sphereqalt2}, each
variable is summed to two constants ($x_1$, $y_1$, $z_1$, $x_2$, $y_2$ and $z_2$). We
need to reduce the number of constants. Let's change variables: instead of $x$, $y$ and $z$,
we are going to use $\alpha$, $\beta$ and $\gamma$ by using the following substitutions:
\begin{equation}\label{eq:sphereqaltsub}
    \begin{cases}
        \alpha = x - x_1\\
        \beta = y - y_1\\
        \gamma = z - z_1
    \end{cases}
    \implies
    \begin{cases}
        x - x_2 = x - x_1 + x_1 - x_2 = \alpha  + (x_1 - x_2)\\
        y - y_2 = y - y_1 + y_1 - y_2 = \beta + (y_1 - y_2)\\
        z - z_2 = z - z_1 + z_1 - z_2 = \gamma + (z_1 - z_2)
    \end{cases}
\end{equation}
This is the difficult part to understand!
In equation \ref{eq:sphereqaltsub}, the substitutions on the left let us replace
the original variables at the numerator of equation \ref{eq:sphereqalt2}; the
calculations on the right are needed so that we can also get rid of the old variables
at the denominator and make the new variables appear:
\begin{equation}\label{eq:sphereqalt3}
    \frac{\alpha^2 + \beta^2 + \gamma^2}{[\alpha  + (x_1 - x_2)]^2 + [\beta + (y_1 - y_2)]^2 + [\gamma + (z_1 - z_2)]^2} = 2
\end{equation}
Now we have a new equation operating on variables $\alpha$, $\beta$ and $\gamma$. Notice that
$x_1 - x_2$, $y_1 - y_2$ and $z_1 - z_2$ are just constants because they are the sum of two
constants. We can introduce another
substitution\sidenote[][-0.29\textwidth]{Notice that we decide to make this
substitution $\Delta_x = x_2 - x_1$ instead of $\Delta_x = x_1 - x_2$ so that, when replacing
in equation \ref{eq:sphereqalt3}, we get $\alpha - \Delta_x$ instead of $\alpha + \Delta_x$.
This will give us a closer result to equation \ref{eq:sphereq}.
However, it is important to point out that both substitutions are fine, even if the final signs will not match
in equation \ref{eq:sphereq}, as long as each variable is summed to a constant that's fine.}:
\begin{equation}\label{eq:sphereqaltsub2}
    \begin{array}{ccc}
        \Delta_x = x_2 - x_1 & \Delta_y = y_2 - y_1 & \Delta_z = z_2 - z_1 \\
    \end{array}
\end{equation}
Which makes equation \ref{eq:sphereqalt3} become:
\begin{equation}\label{eq:sphereqalt4}
    \frac{\alpha^2 + \beta^2 + \gamma^2}{(\alpha  - \Delta_x)^2 + (\beta - \Delta_y)^2 + (\gamma - \Delta_z)^2} = 2
\end{equation}
Clearly, equation \ref{eq:sphereqalt4} has not the same shape as equation \ref{eq:sphereq}.
One thing we need to do for sure, is getting rid of the fraction, which we can
do by multiplying both sides by the denominator:
\begin{equation*}
    \alpha^2 + \beta^2 + \gamma^2 = 2(\alpha  - \Delta_x)^2 + 2(\beta - \Delta_y)^2 + 2(\gamma - \Delta_z)^2
\end{equation*}
Then we can expand each of the 3 squares:
\begin{equation*}
    \alpha^2 + \beta^2 + \gamma^2 =
        2\alpha^2  + 2\Delta^2_x - 4\alpha\Delta_x +
        2\beta^2  + 2\Delta^2_y - 4\beta\Delta_y +
        2\gamma^2  + 2\Delta^2_z - 4\gamma\Delta_z
\end{equation*}
Let's sum up all terms where we can:
\begin{equation}\label{eq:sphereqalt5}
    0 =
        \alpha^2 + \beta^2 + \gamma^2 +
        2\Delta^2_x - 4\alpha\Delta_x +
        2\Delta^2_y - 4\beta\Delta_y +
        2\Delta^2_z - 4\gamma\Delta_z
\end{equation}
Look again at equation \ref{eq:sphereq}, we need to create terms in the
form $(\alpha - c)^2$ ($c$ is just a constant). Looking at equation \ref{eq:sphereqalt5},
we will try now to create\sidenote{To do that, we add and subtract the same quantity.
For example, focusing on the $\alpha$ term only, we make the following quantity:
\begin{equation*}
    \alpha^2 + 2\Delta^2_x - 4\alpha\Delta_x
\end{equation*}
Become:
\begin{equation*}
    \alpha^2 + 2\Delta^2_x - 4\alpha\Delta_x + 2\Delta^2_x - 2\Delta^2_x
\end{equation*}
Which is totally legal. Now, thanks to it, we can create the square:
\begin{equation*}
    (\alpha - 2\Delta_x)^2 - 2\Delta^2_x
\end{equation*}
Which is our aim.} such squared binomals:
\begin{equation*}
    \begin{split}
        0 = &(\alpha^2 + 4\Delta^2_x - 4\alpha\Delta_x) +
            (\beta^2 + 4\Delta^2_y - 4\beta\Delta_y) +
            (\gamma^2 + 4\Delta^2_z - 4\gamma\Delta_z) +\\
            &- 2\Delta^2_x - 2\Delta^2_y - 2\Delta^2_z
    \end{split}
\end{equation*}
We have created the squares, so let's collapse them:
\begin{equation*}
    (\alpha - 2\Delta_x)^2 + (\beta - 2\Delta_y)^2 + (\gamma - 2\Delta_z)^2 =
        2\Delta^2_x + 2\Delta^2_y + 2\Delta^2_z
\end{equation*}
And the equation above is the classic representation of a sphere as in
equation \ref{eq:sphereq}.

\section{Tossing dices}
\label{sec:dices}
\marginnote{Appears as problem 1 in L 06.}
Something

\section{Magnetic fields and amperian loops}
\label{sec:magnamp}
\marginnote{Appears as problem 1 in L 06.}
Something

\section{Spaceships}
\label{sec:spaceship}
\marginnote{Appears as problem 1 in L 06.}
Something

\section{Protons in motion}
\label{sec:protons}
\marginnote{Appears as problem 1 in L 06.}
Something


\chapter[Insights \& curiosity]{Insights \& curiosity}
\label{ch:insights}
There is a lot of Physics in this year's exam, that's why I think we can
take a good chance at delve into some interesting topics.

\section{Continuous bounded functions}
\label{sec:contbound}
Electromagnetism is the pa

\section{Functions symmetric to a generic point}
\label{sec:fsymmcentre}
Electromagnetism is the pa

\section{Polynomial-exponential equations cannot be analytically solved}
\label{sec:polyexpeqimp}
Electromagnetism is the pa

\section{Maxwell's laws explained}
\label{sec:maxwell}
Electromagnetism is the part of Physics that focuses on the study of the
electrical, the magnetic and the electromagnetic fields.


%%%%%%%%%%%%%%%%%%%%%%%%%%%%%%%%%%%%%%%%%%%%%%%%%%%%%%%%%%%%%%%%%%%%%%%%%%%%%%%

\backmatter

\bibliography{biblio}
\bibliographystyle{plainnat}

\printindex

\end{document}
